\documentclass[11pt]{article}
\usepackage[margin=1.1in]{geometry}
\usepackage{amsmath,amssymb,microtype,xcolor,mdframed}
\usepackage[colorlinks=true,linkcolor=blue!50!black,urlcolor=blue!50!black]{hyperref}
\newmdenv[backgroundcolor=blue!4,linecolor=blue!35,linewidth=0.8pt,
  skipabove=9pt,skipbelow=9pt,innerleftmargin=10pt,innerrightmargin=10pt]{keyidea}
\setlength{\emergencystretch}{2em}
\title{Distinct Values Limit Collatz Correction:\\
\Large A Friendly Guide to the Sixth-Root Bound}
\author{Companion to ``A Sixth-Root Bound on Collatz Correction''\thanks{Written with Codex (OpenAI). The technical paper and Lean source give the exact assumptions and proofs. No claim of literature priority is made.}}
\date{September 5, 2026}
\begin{document}
\maketitle
\section{An escaping orbit cannot reuse small numbers}
Use the shortcut Collatz rule: halve an even number; triple an odd number, add one, and halve it. If the sequence ever visits the same number twice, its future repeats. The whole orbit is then bounded.

That means an unbounded orbit must keep visiting distinct numbers. In particular, it cannot repeatedly use the same small odd numbers to collect the effect of the $+1$ in the rule. This simple observation gives a quantitative bound.

Write $N$ for the starting number, $j$ for the number of odd steps taken so far, and $c$ for the normalized affine correction used in the preceding reciprocal-sum paper. The normalization removes the powers of three and two from the accumulated additions.
\begin{keyidea}
\textbf{The new Lean result.} Every unbounded positive orbit obeys
\[
(N+c)^6\le N^6(j+1),
\]
so $c/N$ is at most $(j+1)^{1/6}-1$. No assumption that the reciprocal sum converges is needed.
\end{keyidea}
For illustration, the bound allows $c/N\le1$ after $63$ odd steps, $c/N\le2$ after $728$ odd steps, and $c/N\le9$ after $999{,}999$ odd steps. These are evaluations of the formula, not examples of unbounded Collatz orbits.

The theorem also works on a finite stretch beginning at the seed: if its visited values are distinct and avoid one, the bound holds through the end of that stretch. No assertion about its infinite future is required.

\newpage
\section{Why a sixth root appears}
At an odd value $x$, the quantity $N+c$ is multiplied by
\[
1+\frac1{3x}.
\]
At an even value it is unchanged. For every $x>1$, Lean checks the polynomial inequality
\[
\left(1+\frac1{3x}\right)^6\le\frac{x+1}{x-1}.
\]
For an odd value $x=2k+1$ greater than one, the right side becomes $(k+1)/k$. The odd values are distinct, so the positive integers $k$ are distinct too.

The largest product of $j$ factors of the form $(k+1)/k$ comes from the smallest possible distinct positive $k$ values. Those factors cancel:
\[
\frac21\cdot\frac32\cdot\frac43\cdots\frac{j+1}{j}=j+1.
\]
This bounds the sixth power of the total correction factor by $j+1$. Taking a sixth root gives the announced estimate. The proof uses distinctness and exact arithmetic, not a random-parity model.

\section{What it contributes to the larger proof}
The earlier result said that bounded correction is equivalent to a finite reciprocal orbit sum. It left open the possibility of an escaping orbit whose reciprocal sum diverges. The new theorem applies to that case as well: even if its correction is unbounded, the correction must grow within this slow envelope.

There is also a constraint on odd steps. After $L$ steps, if an unbounded orbit has not dropped below its starting value, then
\[
2^{6L}\le3^{6j}(j+1).
\]
This places the odd-step count close to or above the familiar critical line, allowing only a logarithmic deficit. Lean proves the displayed integer inequality directly.
\begin{keyidea}
\textbf{What is still missing.} Slow growth is still growth. The sixth-root bound does not make correction bounded, and it does not force descent. We still need an arithmetic argument that excludes the itineraries allowed by these estimates. Nontrivial cycles remain a separate question.
\end{keyidea}
The implementation is \texttt{lean/Collatz/CorrectionGrowth.lean}. The technical paper contains the complete proof and rebuild instructions. The expanded audit covers 53 selected declarations with only standard foundational axioms. Collatz has not been proved.
\end{document}
